\documentclass[parskip=full]{scrartcl}
\usepackage[top=2.5cm, bottom=2.5cm, left=2.5cm, right=2.5cm]{geometry}
\usepackage[utf8]{inputenc}
\usepackage[T1]{fontenc}
\usepackage[german]{babel}
\usepackage{hyperref}
\usepackage[toc, nonumberlist]{glossaries}
\usepackage{graphicx}
\usepackage{enumitem}
\usepackage{float}
\usepackage{color}
\usepackage{textcomp}

\hypersetup{
  colorlinks=false,
  linktoc=all,
  hidelinks,
}

\title{OSIP - Testbericht}
\subtitle{OPC UA Simulator for Industrial Plants}
\author{
    M. Armbruster\\
    D. Kahles\\
    H. Lehmann\\
    M. Schwarzmann\\
    N. Wilhelm
}

\begin{document}
\maketitle
\thispagestyle{empty}
\vspace{20px}
\begin{center}
  \includegraphics[scale=0.4]{../icon.png}
\end{center}
\pagebreak
\tableofcontents
\pagebreak

\section{Einleitung}
Dieses Dokument behandelt die Qualitätssicherung von OSIP.
Das Dokument umfasst Änderungen in der Implementierung, die in der Testphase vorgenommen wurden, sowie Unittests.
Im Übrigen werden die betrachteten Anforderungen und geprüften Testfälle erläutert.
Insbesondere wird aufgezeigt, an welchen Stellen sich die Anforderungen und Testfälle im Vergleich zum Pflichtenheft verändert haben. Einige Anforderungen und Testfälle wurden überarbeitet oder erweitert, damit sie besser zur Implementierung passen.
Das Dokument beschreibt die geprüften Szenarien, deren Ausgang und die Abdeckung des Codes bei der Durchführung der Szenarien.
Im Anschluss befindet sich ein Fehlerprotokoll, wo aufgetretene Fehler, ihre Ursache und ihre Behandlung dokumentiert sind.
\section{Änderungen in der Implementierung}
Während der Testphase wurden einige Fehler behoben und Verhaltensweisen verändert, um OSIP benutzerfreundlicher zu machen.
Diese Probleme wurden nicht notwendigerweise durch die Testszenarien aufgedeckt, da die Verhaltensweisen nicht im Fokus des
Pflichtenhefts standen und daher nicht so genau spezifiziert wurden (u.a. da Fragen zu Details der GUI oder zum Verhalten bei Netzwerkproblemen
erst im Laufe der Entwicklung aufkamen) oder die Szenarien die Fehler nicht aufdecken konnten, weil sie nur durch bestimmte,
ungewöhnliche Eingaben ausgelöst werden konnten. Die wichtigsten Änderungen sind die folgenden:

\begin{itemize}
 \item Vereinfachen von Übersetzungen von OSIP in andere Sprachen, indem die lokalisierten Fehlermeldungen sämtlicher Exceptions benutzt werden.
 \item Beheben von Abstürzen, z.b. wenn invalide Werte in den Einstellungsdateien stehen.
 \item Öffnen sämtlicher Dialoge (wie z.B. Alarmmeldungen, Einstellungen oder Überlaufmeldungen) über den anderen Fenstern und modale Initialisierung dieser Dialoge,
 sodass OSIP nicht bedienbar ist, solange diese nicht geschlossen werden. Das verhindert beispielsweise, dass der Benutzer nach einem Verbindungsverlust
 die Einstellungen nicht bedienen kann, weil der Fehlerdialog im Hintergrund unbemerkt offen ist.
 \item Verbesserte Logik beim Starten der OPC UA Server und Clients. Wenn ein Server Port mehrfach verwendet wird, wird sofort das
 Einstellungsfenster geöffnet.
 \item Performanceverbesserungen bei den OPC UA Servern und Clients und beim Beenden von OSIP.
 \item Setzen des OSIP Icons in den Fensterrahmen jedes Fensters.
 \item Verbessern der mitgelieferten Szenarien.
 \item Besseres Visualisieren von ausgeschalteten Verläufen.
 \item Vereinfachen des Codes an einigen Stellen.
\end{itemize}

\section{Unittests}
Die Phase der Qualitätssicherung haben wir damit begonnen, die Unittests aus der Implementierungsphase weiter auszubauen und so die Testabdeckung zu maximieren.
Innerhalb des Models war die Testabdeckung bereits sehr gut. Für die Controller und die graphischen Benutzeroberflächen waren noch keine Unittests vorhanden.

Für die Controller blieben wir bei der Entscheidung, keine Unittests einzuführen, da diese hier kaum anwendbar waren. Stattdessen schrieben wir für die Controller
Integrationstests. Für die beiden graphischen Benutzeroberflächen haben wir jedoch noch Testklassen
erstellt. Hier konnte nur eine geringere Abdeckung erreicht werden, als es im Model der Fall ist, dennoch liegt in beiden Oberflächen die Zeilenabdeckung bei
über 50\%.

Die endgültige Zeilenabdeckung der Module, für die Unittests verwendet wurden, ist in der folgenden Tabelle nochmals dargestellt. Dabei werden sowohl die
Werte vor als auch die Werte nach der Testphase gelistet.

\begin{tabular} {| l || r | r |}
  \hline
  Modul				& Vorher	& Nachher	\\
  \hline \hline
  osip-core			& 82.6\%	& 86.1\%	\\
  osip-model			& 91.8\%	& 96.6\%	\\
  osip-model-Simulation		& 100.0\%	& 100.0\%	\\
  osip-monitoring-view		& 0.0\%		& 82.3\%	\\
  osip-simulation-controller	& 0.0\%		& 6.3\%		\\
  osip-simulation-view		& 0.0\%		& 50.2\%	\\
  \hline
\end{tabular}

\section{Anforderungen und Tests}
In Folge unterschiedlicher Umsetzungen von Anforderungen im Vergleich zum Pflichtenheft erfolgt zunächst die Anpassung der Anforderungen und Tests durch die vorhandene Rückkoppelung.

\subsection{Angepasste Anforderungen}
Die Anforderungen [FA100], [FA110], [FA310], [FA320], [FA450], [FA610], [FA720], [FA730], [FA740], [NF100], [D20], [D60], [D70], [D100] und [D120], die zu den Kannkriterien zählten, wurden nicht implementiert und konnten daher gestrichen werden.

Folgende Anforderungen wurden umgeändert (Änderungen bzw. Hinzufügungen sind \textbf{fett} markiert):
\begin{enumerate}
	\item[FA290] Der Menüeintrag "`Szenarien"' \textbf{öffnet einen Dateidialog zur Auswahl eines Simulations-Szenarios, das anschließend gestartet wird.}
	\item[FA560] Die Überwachungskonsole registriert für jeden Tank jeweils einen Alarm für einen Überlauf und für einen Unterlauf \textbf{sowie einen für Überhitzung und Unterkühlung}.
	\item[FA580] Ein Alarm wird durch einen Bezeichner und einen Kreis als Zustandssymbol dargestellt. Ist der Kreis rot, wurde der Alarm ausgelöst, ist er grün, wurde der Alarm nicht ausgelöst. \textbf{Ist er grau, ist der Alarm ausgeschalten.}
	\item[FA630] Ein Alarm wird als ausgelöst angezeigt, nachdem die Überwachungskonsole durch die Fertigungssimulation über den Alarm informiert wurde. Die Anzeige besteht, bis der Alarm in der Fertigungssimulation nicht mehr aktiv ist und \textbf{die Überwachungskonsole über das Ende informiert wurde}.
	\item[FA640] Es gibt eine zweistufige Ampel für den allgemeinen Zustand der Überwachungskonsole. Ist die Ampel rot, wurde mindestens ein Alarm ausgelöst, ist sie grün, wurde kein Alarm ausgelöst. \textbf{Ist sie grau, sind alle Alarme ausgeschalten.}
	\item[FA700] In den Einstellungen wird die IP-Adresse der OPC UA Server sowie deren Ports angezeigt. Die Ports und die IP-Adresse der OPC UA Server sind konfigurierbar.
	\item[FA780] Der Nutzer kann die verfügbaren Alarme ein- und ausschalten.
	\item[FA800] Der Benutzer wird über eine Fehleingabe informiert \textbf{oder es erfolgt eine automatische Korrektur in den Wertebereich}.
	\item[NF40] Standardmäßig haben die OPC UA Server die Ports \textbf{12868, 12869, 12870 und 12871}.
	\item[NF60] Die Drehzahl des Mischermotors kann vom Benutzer in einem Bereich zwischen 0 und \textbf{3000} Umdrehungen pro Minute gewählt werden.
	\item[NF90] Der Zu- und Abfluss der oberen Tanks und der Abfluss des unteren Tankes liegen zwischen 0 und \textbf{25} Litern pro Sekunde. Der Zufluss des unteren Tanks liegt zwischen 0 und \textbf{75} Litern pro Sekunde.
	\item[NF110] Die Schwellenwerte für die Füllstands-Alarme liegen zwischen 0 und 100 \% \textbf{und für die Temperatur-Alarme innerhalb von 0 \textcelsius und 230 \textcelsius}. Standardmäßig liegt der Schwellenwert für einen Überlauf bei 95 \%, für einen Unterlauf bei 5 \%\textbf{, für die Überhitzung bei 95 \% und für die Unterkühlung bei 5 \% des einstellbaren Temperaturintervalls}.
	\item[NF140] Die Aktualisierungsfrequenz kann vom Benutzer im Bereich zwischen \textbf{500} und 4000 ms gewählt werden.
	\item[D40] Die Überwachungskonsole speichert die zuletzt eingeschalteten Alarme.
\end{enumerate}

\subsection{Angepasste globale Testfälle}
Die globalen Testfälle [GT160], [GT170], [GT200] und [GT240] wurden aus den Testszenarien gestrichen und nicht durchgeführt, da sie auf nicht implementierten Kannkriterien basieren.
Zudem wurden alle nicht umgesetzten Anforderungen in den globalen Testfällen nicht überprüft. Explizit wurden folgende Testfälle leicht umgeändert (Änderungen auch hier \textbf{fett}):

\begin{enumerate}
	\item[GT210] Korrekter Ablauf der Simulations-Szenarien: Die Fertigungssimulation läuft und die Überwachungskonsole ist bei dieser registriert. \textbf{Alle verfügbaren, d. h. von uns mit gelieferten, Simulations-Szenarien} werden nacheinander korrekt ausgeführt (FA120, FA290), ohne dass \textbf{während der Ausführung eines Szenarios} eine Bedienung durch den Nutzer benötigt wird.
	\item[GT270] Fehlerhafte Eingaben: Alle möglichen Einstellungen in der Fertigungssimulation und Überwachungskonsole werden mit Werten außerhalb des Wertebereichs vorgenommen (NF10, NF20, NF60, NF90, NF110, NF140). Jedes Mal muss eine Fehlermeldung erscheinen (FA790, FA800) \textbf{oder eine automatische Korrektur in den Wertebreich erfolgen (FA800)}.
	\item[\textbf{GT290}] \textbf{Dauerhaftes Speichern der Einstellungen: Die Fertigungssimulation und Überwachungskonsole werden mit Standardeinstellungen geöffnet. Alle Einstellungen werden geändert und gespeichert. Nach einem Neustart beider Programme sind die geänderten Einstellungen weiterhin vorhanden (D10, D30 - D50, D110, D130).}
\end{enumerate}

\subsection{Testszenarien}
Im Nachfolgenden werden eventuelle Änderungen der Testszenarien sowie das Ergebnis dieser dargestellt. Alle Szenarien wurden unter Ubuntu 16.04 und Windows 10 durchgeführt.

\subsubsection{Szenario 1: Fertigungssimulation}
\begin{description}
	\item[Ablauf] [GT100], [GT280].
	\item[Ergebnis] Erfolgreich.
	\item[Abdeckung] 15,7 \% der Zeilen.
\end{description}

\subsubsection{Szenario 2: Überwachungskonsole mit ungültigen Einstellungen}
\begin{description}
	\item[Ablauf] [GT100], [GT120], [GT280].
	\item[Ergebnis] Erfolgreich.
	\item[Abdeckung] 54,8 \% der Zeilen.
\end{description}

\subsubsection{Szenario 3: GUI}
\begin{description}
	\item[Ablauf] [GT100], [GT110], [GT140], [GT280].
	\item[Ergebnis] Legte den Fehler HL-003 offen.
	\item[Abdeckung] 69,6 \% der Zeilen.
\end{description}

\subsubsection{Szenario 4: Verbindungsinformationen}
\begin{description}
	\item[Ablauf] [GT100], [GT110], [GT150], [GT260], [GT280].
	\item[Ergebnis] Legte den Fehler MA-002 offen.
	\item[Abdeckung] 46,1 \% der Zeilen.
\end{description}

\subsubsection{Szenario 5: Sensordaten}
\begin{description}
	\item[Ablauf] [GT100], [GT110], [GT180], [GT280].
	\item[Ergebnis] Erfolgreich.
	\item[Abdeckung] 68,6 \% der Zeilen.
\end{description}

\subsubsection{Szenario 6: Alarme und Logging}
\begin{description}
	\item[Ablauf] [GT100], [GT110], [GT190], [GT230], [GT260], [GT280].
	\item[Ergebnis] Erfolgreich.
	\item[Abdeckung] 59 \% der Zeilen.
\end{description}

\subsubsection{Szenario 7: Aktualisierungsfrequenz und Verläufe}
\begin{description}
	\item[Ablauf] [GT100], [GT110], [GT220], [GT250], [GT280].
	\item[Ergebnis] Legte den Fehler DK-001 offen.
	\item[Abdeckung] 60 \% der Zeilen.
\end{description}

\subsubsection{Szenario 8: Fehlerhafte Eingaben}
\begin{description}
	\item[Ablauf] [GT100], [GT110], [GT270], [GT280].
	\item[Ergebnis] Erfolgreich.
	\item[Abdeckung] 58,2 \% der Zeilen.
\end{description}

\subsubsection{Szenario 9: Simulations-Szenarien}
\begin{description}
	\item[Ablauf] [GT100], [GT110], [GT210], [GT280].
	\item[Ergebnis] Erfolgreich.
	\item[Abdeckung] 81,3 \% der Zeilen.
\end{description}

\subsubsection{Szenario 10: Dockerimage}
\begin{description}
	\item{Anmkerung} Dem Pflichtenheft entsprechend, nur unter Ubuntu 16.04 getestet.
	\item[Ablauf] [GT130], [GT280].
	\item[Ergebnis] Erfolgreich.
	\item[Abdeckung] 40,7 \% der Zeilen.
\end{description}

\subsubsection{Szenario 11: Produktdaten}
\begin{description}
	\item[Anmerkung] Neues Szenario, um die Produktdaten zu testen. Diese wurden im Pflichtenheft vernachlässigt.
	\item[Ablauf] [GT100], [GT110], [GT290], [GT280].
	\item[Ergebnis] Erfolgreich.
	\item[Abdeckung] 71,7 \% der Zeilen.
\end{description}

\section{Fehlerprotokoll}

\subsection{DK-001}
\begin{description}
	\item[Beschreibung] [GT220] schlägt fehl, da die Verzögerung zu groß ist.
	\item[Ursache] In Milo wird das eingestellte Intervall beim Subscriben an zwei Stellen verwendet. Dadurch kann es dazu kommen, dass die Verzögerung bis zu doppelt so groß ist als eingestellt.
	\item[Behebung] Das Milo-interne Intervall, mit dem die Änderungen im Modell erfasst werden, wird auf konstant 250 Millisekunden gesetzt.
\end{description}

\subsection{MA-001}
\begin{description}
	\item[Beschreibung] Der Unittest edu.kit.pse.osip.io.files.ScenarioFileTest schlägt unter Windows fehl.
	\item[Ursache] Der Pfad der Testdatei ist nicht für Windows konfiguriert.
	\item[Behebung] Der Pfad wird in eine java.io.File-Instanz konvertiert, die einen passenden, absoluten Pfad für jedes Betriebssystem liefert.
\end{description}

\subsection{MA-002}
\begin{description}
	\item[Beschreibung] Nach Einstellung neuer Ports für die OPC UA Server in der Fertigungssimulation und fertigem Neustart der OPC UA Server können sich die OPC UA Clients nach Einstellen der neuen Ports und deren Speichern in der Überwachungskonsole nicht verbinden. Auch nach dem erneuten Speichern erfolgt keine Herstellung der Verbindung.
	\item[Ursache] Die Spinner, mit denen die Ports eingestellt werden, übernehmen intern nicht den eingegebenen Wert, wenn sie den Fokus verlieren.
	\item[Behebung] Ein zusätzliches bidirektionales Binding zwischen dem TextFormatter und der ValueFactory wurde hinzugefügt.
\end{description}

\subsection{HL-001}
\begin{description}
	\item[Beschreibung] Das UI von Überwachungskonsole und Simulation ist nach dem Start für ca. 3-10 Sekunden blockiert und reagiert nicht. Nach dem Speichern der Einstellungen tritt es ebenfalls auf.
	\item[Ursache] Die OPC UA Server und Clients werden im UI Thread gestartet und beendet.
	\item[Behebung] Starte einen neuen Thread für Aktionen mit OPC UA. Zeige während der Aktionen einen Warte-Dialog an, um den Nutzer zu informieren.
\end{description}

\subsection{HL-002}
\begin{description}
	\item[Beschreibung] Während die Einstellungen der Simulation angewendet werden, bleibt das Einstellungsfenster offen und kann modifiziert werden.
	\item[Ursache] Das Einstellungsfenster wird nicht geschlossen.
	\item[Behebung] Einstellungsfenster schließen und bei fehlerhaften Einstellungen wieder öffnen.
\end{description}

\subsection{HL-003}
\begin{description}
	\item[Beschreibung] Beim Schließen der Einstellungen der Simulation werden die Werte nicht zurückgesetzt, erneutes Öffnen zeigt die alten Werte an.
	\item[Ursache] Die Werte werden nicht zurückgesetzt.
	\item[Behebung] Vor dem Anzeigen des Fensters werden die Werte zurückgesetzt.
\end{description}

\subsection{HL-004}
\begin{description}
	\item[Beschreibung] Einige Szenarien-Dateien können nicht gelesen werden.
	\item[Ursache] Die Dateien enthalten eine UTF-8 ByteOrderMark, die der Parser nicht erwartet.
	\item[Behebung] Beginnt eine Datei mit einer ByteOrderMark, so werden die Zeichen beim Parsen übersprungen.
\end{description}

\pagebreak
\phantomsection
\addcontentsline{toc}{section}{\listfigurename}
\listoffigures

\end{document}
